% Nouveau EGA V (cote 25): archive cote 25, site page 166 (PDF page 167).
% typescript. Draft compared: none.
% AI-transcribed by Claude Opus 5.5 (claude-opus-5-5, adaptive thinking, effort high), Message Batches API, batch ega25b (msgbatch_01A7n5YuXubUPNK1zsepi8RG), 2026-11-13; original AI transcription, subsequently checked against both supplied scan parts; see PAGES_CHECKED.tsv and CHECK_DECISIONS.tsv.
\scanpage{25/166}
\placedmargin{\unsure{//}} \hand{nouveau EGA $\overline{\underline{\mathrm{V}}}$} \mspage{\hand{1}}

\placedmargin{\unsure{(}}EGA IV par ?

\uline{Ensemble singulier et supersingulier}.

Pages 14 et 15 des notes précédentes, remplacer partout S par Y. Page
15 ligne 10 le \hand{``}i.e.\hand{''} est ab\unsure{u|e}\unsure{s|u}\unsure{i|f}f, c'est la deuxième assertion qui est la
plus précise et celle à laquelle on en a.

La généralisation de prop.7 à laquelle il est fait allusion pages 16
et 17 s'obtient à l'aide du diagramme suivant, qui généralise celui de la
page 4 :
\diagram{\begin{tikzcd}[ampersand replacement=\&,nodes in empty cells]
 \& \& \hand{0} \& \& \\
 \& \& \hand{\underline{\Omega}}^1_{X/Y}\hand{\otimes}\underline{O}_{V^1} \arrow[u,handwritten] \& \& \\
0 \arrow[r,handwritten] \& \underline{N}_{V^1/X} \arrow[r,handwritten] \arrow[ur,handwritten, "\hand{\mu}"] \& \hand{\underline{\Omega}}^1_{X/S}\hand{\otimes}\underline{O}_{V^1} \arrow[r,handwritten] \arrow[u,handwritten] \& \hand{\underline{\Omega}}^1_{V^1/S} \arrow[r,handwritten] \& 0 \\
 \& \& \hand{\underline{\Omega}}^1_{Y/S}\hand{\otimes}\underline{O}_{V^1} \arrow[u,handwritten] \arrow[ur,handwritten, "\hand{\wr}"'] \& \& \\
 \& \& \hand{0} \arrow[u,handwritten] \& \&
\end{tikzcd}}

où $\underline{N}$ désigne le faisceau conormal. Tous les faisceaux du diagramme sont
localement libres sur $V^1$, et on app\unsure{l|D}ique le lemme de la page 7 à $M=\hand{\underline{\Omega}}^1_{X/S}\hand{\otimes}\underline{O}_{V^1}$
et aux deux sous-Modules figurant dans le diagramme précédent. Avec
les notations de ce lemme, on est ramené à prouver que la section a
 de \uline{Hom}(L,P) déduit de $\unsure{\mu|u}$: $P\hand{\rightarrow} M/Q$ i.e. de $\mu$: $\underline{N}_{V^1/X}\hand{\longrightarrow}\hand{\underline{\Omega}}^1_{X/Y}\hand{\otimes}\underline{O}_{V^1}$ ,
dont le préschéma des zéros V(a) est celui qui rend \hand{$\mu$} strictement de rang
\hand{$<$} d = rang M/Q , satisfait V(a) = V(D), où D est défini au bas de la
page 15. Or les inclusions de SS-préschémas lisses
\placedmargin{||}
\[ V^1 \hand{\subset} V^o \hand{\subset} X \]
donnent une suite exacte (réf ?)
\[ 0 \hand{\rightarrow} \underline{N}_{V^o/X}\hand{\otimes}\underline{O}_{V^1} \hand{\rightarrow} \underline{N}_{V^1/X} \hand{\longrightarrow} \underline{N}_{V^1/V^o} \hand{\rightarrow} 0 \ , \]
d'autre part on a des isomorphismes canoniques
\[ \underline{N}_{V^o/X} \hand{\simeq} \underline{E}\hand{\otimes}\underline{O}_{V^1} \ , \quad \underline{N}_{V^1/V^o} \hand{\simeq} \uline{Hom}(W,K) \ , \]
où W,K sont définis (page 15) comme noyau et conoyau de

\mspage{\hand{166}}
